\documentclass[../main.tex]{subfiles}

\begin{document}
\begin{CJK*}{UTF8}{gbsn}

\section*{Exercise 7}
Consider the conformal prediction with $\mathcal{Y} = \mathbb{R}$
with a given model $\hat{f}: \mathcal{X} \to \mathcal{Y}$
and calibration points $(X_1, Y_1)$,$\ldots$, $(X_n, Y_n)$.
Put 

\begin{equation*}
    \hat{q}_L = \text{Quantile}\bracket{\curly{\hat{f}(X_i)-Y_i}_{i=1}^n ; 
    \bracket{1- \frac{\alpha}{2}} \bracket{1+ \frac{1}{n}}},
\end{equation*}

\begin{equation*}
    \hat{q}_R = \text{Quantile}\bracket{Y_i-\curly{\hat{f}(X_i)}_{i=1}^n ; 
    \bracket{1- \frac{\alpha}{2}} \bracket{1+ \frac{1}{n}}},
\end{equation*}
and
\begin{equation*}
    \mathcal{C}(X_{n+1}) = \squared{\hat{f}(X_{n+1}) - \hat{q}_L, \hat{f}(X_{n+1}) + \hat{q}_R}.
\end{equation*}
Prove that 

\begin{equation*}
    P \curly{Y_{n+1} \in \mathcal{C}(X_{n+1})} \geqslant 1-\alpha
\end{equation*}
by showing both

\begin{equation*}
    P \curly{Y_{n+1} < \hat{f}(X_{n+1}) - \hat{q}_L} \leqslant \frac{\alpha}{2}
\end{equation*}
and 
\begin{equation*}
    P \curly{Y_{n+1} > \hat{f}(X_{n+1}) + \hat{q}_R} \leqslant \frac{\alpha}{2},
\end{equation*}
assuming that $(X_1, Y_1),\ldots, (X_{n+1}, Y_{n+1})$
is exchangeable.

\paragraph{Solution}
We have

\begin{equation*}
    P \curly{Y_{n+1} < \hat{f}(X_{n+1}) - \hat{q}_L}
    = P \curly{\hat{f}(X_{n+1}) - Y_{n+1} > \hat{q}_L }
    = 1 - P \curly{\hat{f}(X_{n+1}) - Y_{n+1} \leqslant \hat{q}_L }
    \leqslant 1 - \bracket{1- \frac{\alpha}{2}} = \frac{\alpha}{2}
\end{equation*}
where the last inequality 
uses Theorem 3.2 in the textbook
with score $S_i = \hat{f}(X_{i}) - Y_{i}$.
This is legitimate because the data is assumed 
to be exchangeable and the score is trivially 
symmetric in case of split conformal prediction.
The right hand side inequality is proved in 
exactly the same way but with score $S_i = Y_i-\hat{f}(X_{i})$.
Combine the two, we have the desired marginal coverage.
The proof is complete.
////
\end{CJK*}
\end{document}